\section{Physical Basis for Nuclear-State Control of Fission}
\label{sec:physical_basis}

The framework requires two independent physical ingredients. First, nuclear states must be externally addressable or controllable in at least some systems. Second, fission probability must depend on the populated nuclear state. Neither statement is intrinsically speculative; what remains unestablished is whether the two can be combined in a regime that yields useful, externally controlled modification of fission.

\subsection{External manipulation of nuclear states}

The first requirement is
\begin{equation}
u_Q(t)
\longrightarrow
\delta\rho_N(t).
\end{equation}

Nuclear quantum optics provides direct evidence that nuclear excitation dynamics can be manipulated by external fields. Bürvenich, Evers, and Keitel showed theoretically that high-frequency laser facilities could enable resonant laser--nucleus interactions and explicitly identified preparation and control of nuclear states as a long-term objective \cite{Buervenich2006}.

Experimentally, coherent X-ray-optical control of collective nuclear excitations has been demonstrated in $^{57}$Fe. Heeg \textit{et al.} used phase-controlled X-ray pulses to switch nuclear dynamics between enhanced coherent excitation and stimulated emission \cite{Heeg2021}. This establishes that phase-sensitive manipulation of a nuclear excitation is experimentally possible.

The $^{229}$Th system provides a complementary result. Its unusually low-energy nuclear isomeric transition lies in the optical/VUV domain, and high-resolution laser spectroscopy of the isomeric transition has now been achieved \cite{Yamaguchi2024}. Although $^{229}$Th is not itself the principal fissile candidate considered below, it demonstrates that the internal state of a nucleus is not universally inaccessible to external optical control.

At much higher excitation energies, Pálffy \textit{et al.} investigated coherent MeV photon pulses interacting with heavy nuclei in a quasiadiabatic regime where photon absorption competes with intrinsic nuclear equilibration. Their model includes neutron evaporation and fission among the decay channels \cite{Palffy2015}. This work does not demonstrate controlled fission, but it shows that externally driven coherent excitation and fission can occur within the same dynamical framework.

\subsection{Fission probability is state dependent}

The second requirement is
\begin{equation}
\frac{\partial P_f}{\partial\rho_N}
\neq0.
\end{equation}

In compound-nucleus reaction theory, decay probabilities depend not only on excitation energy but also on quantities such as angular momentum, parity, transition-state structure, and available channels. Schematically,
\begin{equation}
\sigma_{n,f}(E)
=
\sum_{J,\pi}
\sigma^{CN}(E,J,\pi)
G_f(E^\ast,J,\pi),
\end{equation}
where $\sigma^{CN}$ denotes compound-nucleus formation and $G_f$ the probability of decay through fission.

Thus, even in a statistical description,
\begin{equation}
G_f(E^\ast,J_1,\pi_1)
\neq
G_f(E^\ast,J_2,\pi_2)
\end{equation}
in general.

The strongest anchor for the present work is more direct: the low-lying isomer of $^{235}$U has a measured neutron-induced fission cross section different from that of the ground state \cite{DEer1988}. Consequently, the link
\begin{equation}
\rho_N\rightarrow\sigma_f
\end{equation}
is experimentally established in a relevant fissile isotope.

\subsection{Fission as quantum many-body dynamics}

Modern microscopic theory treats induced fission as large-amplitude, time-dependent many-body dynamics \cite{SchunckRobledo2016}. Time-dependent nuclear density-functional calculations follow the evolution toward scission and show sensitivity to the initial configuration and dynamical trajectory. Finite-temperature TDDFT studies, for example, explicitly investigate dissipation and trajectory-dependent evolution during induced fission \cite{Li2023}.

Alternative microscopic descriptions also emphasize that only subsets of many-body states may contribute strongly to transmission across the fission barrier. A nonequilibrium Green-function treatment of $^{236}$U found that the fluctuations of fission probability can be associated with the number of many-body eigenstates that simultaneously have significant support on both sides of the barrier and lie near the system excitation energy \cite{UzawaHagino2024}.

These results motivate the search for pairs or families of configurations satisfying
\begin{equation}
P_f^{(A)}
\neq
P_f^{(B)}
\end{equation}
together with a physically accessible external transition or redistribution
\begin{equation}
|\Psi_A\rangle
\overset{u_Q}{\longleftrightarrow}
|\Psi_B\rangle.
\end{equation}

\subsection{Population control versus coherent control}

Two different physical regimes should be distinguished.

In the weaker, population-based regime, the density matrix is effectively diagonal,
\begin{equation}
\rho_N
=
\sum_i
p_i
|i\rangle\langle i|,
\end{equation}
and
\begin{equation}
P_f^{\mathrm{eff}}
=
\sum_i
p_i P_f^{(i)}.
\end{equation}

Here, genuine long-lived coherence is unnecessary. Redistribution of state populations is sufficient if the states have different fission responses.

In the stronger regime, the external interaction creates a coherent superposition
\begin{equation}
|\Psi\rangle
=
\sum_i c_i|i\rangle
\end{equation}
with nonzero off-diagonal density-matrix elements. If multiple components feed common fission channels, the transition amplitude can contain interference terms. Schematically,
\begin{equation}
A_f
=
\sum_i c_i A_{i\rightarrow f},
\end{equation}
so that
\begin{equation}
P_f
=
|A_f|^2
\end{equation}
contains cross terms depending on relative phases.

A phase-dependent fission response would represent genuine coherent control, but it is also the more demanding scenario because highly excited heavy nuclei possess dense spectra and rapid internal equilibration. The overall control concept therefore does not require coherent macroscopic superpositions: population control alone is sufficient to define an additional state-dependent fission-control channel.

\subsection{The missing connection}

Existing physics supports the statements
\begin{equation}
\boxed{\text{nuclear states can be externally manipulated}}
\end{equation}
and
\begin{equation}
\boxed{\text{fission response can depend on nuclear state}}.
\end{equation}

What remains to be established in a suitable fissile system is the controlled connection
\begin{equation}
\boxed{
\begin{minipage}{0.86\columnwidth}
\centering
external actuation $\rightarrow$ controlled nuclear-state redistribution\\
$\rightarrow$ controlled modification of neutron-induced fission
\end{minipage}
}.
\end{equation}

This missing connection defines the central microscopic problem of the present framework.
